\newgeometry{top=1.5cm, bottom=1.5cm, left=3cm, right=3cm} 
\chapter{Conclusioni}
\begin{flushleft}
    La presente tesi ha fornito una panoramica dettagliata sull'applicazione degli algoritmi genetici per l'addestramento delle reti neurali. Sono stati analizzati approfonditamente i principi teorici del machine learning, le architetture delle reti neurali e i meccanismi degli algoritmi evolutivi, presentando inoltre gli strumenti software utilizzati per lo sviluppo del progetto e la struttura metodologica del lavoro svolto.
    L'obiettivo principale di questo lavoro non è stato quello di sviluppare una soluzione  superiore alle tecniche esistenti, ma piuttosto di esplorare e validare un approccio alternativo all'addestramento delle reti neurali. La ricerca si è concentrata nel dimostrare come gli algoritmi genetici possano essere  utilizzati per l'ottimizzazione dei pesi delle reti neurali, offrendo un metodo alternativo rispetto alle consolidate tecniche di backpropagation tradizionali.\\
    Per garantire un confronto metodologicamente coerente tra i diversi approcci, è stata  utilizzata la funzione di cross-entropy come metrica di valutazione, pur non sfruttando completamente le potenzialità degli algoritmi genetici. Va sottolineato infatti che la caratteristica peculiare e distintiva di questi algoritmi è la loro capacità di operare efficacemente senza richiedere la differenziabilità della funzione obiettivo, proprietà particolarmente vantaggiosa in numerosi contesti applicativi come il reinforcement learning, dove questa caratteristica matematica non è sempre disponibile o facilmente implementabile.\\
    I risultati sperimentali ottenuti hanno evidenziato alcuni aspetti interessanti degli algoritmi genetici nell'addestramento delle reti neurali, sebbene nel caso studio specifico analizzato non siano riusciti a raggiungere le prestazioni del metodo tradizionale di riferimento. Tuttavia, è emerso come la combinazione di entrambi gli approcci metodologici possa potenzialmente portare a risultati più promettenti rispetto all'utilizzo esclusivo di una singola tecnica.
\end{flushleft}